% =========================================================
%  Presentación de avances — Proyecto SIA (Biblioteca)
%  Compilar con: pdflatex avances_sia.tex  (dos veces)
%  Para los diagramas: exportar desde Bizagi / la herramienta
%  del DFD como PNG o PDF y dejarlos en la carpeta figuras/
%  con los nombres indicados. Si no existen, se muestra un
%  recuadro de reemplazo.
% =========================================================
\documentclass[aspectratio=169,10pt]{beamer}

\usepackage[spanish,es-noshorthands]{babel}
\usepackage[utf8]{inputenc}
\usepackage[T1]{fontenc}
\usepackage{lmodern}
\usepackage{graphicx}
\usepackage{booktabs}
\usepackage{tabularx}
\usepackage{tikz}

\usetheme{default}
\usecolortheme{default}
\setbeamertemplate{navigation symbols}{}
\setbeamertemplate{footline}[frame number]

\definecolor{azul}{RGB}{0,61,121}
\definecolor{naranjo}{RGB}{234,114,0}
\definecolor{gris}{RGB}{110,110,110}
\setbeamercolor{frametitle}{fg=azul}
\setbeamercolor{title}{fg=azul}
\setbeamercolor{structure}{fg=azul}
\setbeamercolor{block title}{bg=azul,fg=white}
\setbeamercolor{block body}{bg=azul!6}
\setbeamerfont{frametitle}{series=\bfseries}

% Figura o recuadro de reemplazo si el archivo aún no existe
\newcommand{\figura}[3]{%
  \IfFileExists{#1}{%
    \includegraphics[width=#2,height=#3,keepaspectratio]{#1}%
  }{%
    \begin{tikzpicture}
      \draw[dashed,gris,thick,rounded corners] (0,0) rectangle (#2,#3);
      \node[gris,align=center] at (0.5*#2,0.5*#3)
        {\small Insertar diagrama\\\footnotesize\texttt{\detokenize{#1}}};
    \end{tikzpicture}%
  }%
}

\title{Diagnóstico de Sistemas de Información}
\subtitle{Avance del proyecto: Biblioteca}
\author{\small Integrantes: Iván Mansilla, Ayrton Morrison}
\institute{Diseño de Sistemas de Información}
\date{Octubre 2026}

\begin{document}

% ---------------------------------------------------------
\begin{frame}
  \titlepage
\end{frame}

% ---------------------------------------------------------
\begin{frame}{Contenido}
  \tableofcontents
\end{frame}

% =========================================================
\section{Descripción de la empresa}
% =========================================================
\begin{frame}{Descripción de la organización}
  \begin{columns}[T]
    \begin{column}{0.55\textwidth}
      \textbf{Biblioteca universitaria}\\
      \small Unidad de Gestión de Recursos Académicos

      \vspace{0.6em}
      \begin{itemize}\small
        \item Gestiona el recurso bibliográfico físico y digital de la universidad.
        \item Presta libros, espacios (logias, puestos de estudio) y equipos
              (notebooks, bicicletas, audífonos).
        \item Apoya a las carreras en acreditación y en la compra de bibliografía.
        \item Cumple un rol social: actividades abiertas a la comunidad.
      \end{itemize}
    \end{column}
    \begin{column}{0.42\textwidth}
      \begin{block}{Sistema principal}
        \small
        \textbf{Horizont}: catálogo MARC21 (Dewey/Cutter), ítems, usuarios,
        préstamos, multas y estados del ítem.
      \end{block}
      \vspace{0.3em}
      \begin{block}{Apoyo informal}
        \small Planillas Excel y Drive para registro de ingreso, inventario,
        actas y libros digitalizados.
      \end{block}
    \end{column}
  \end{columns}
\end{frame}

% =========================================================
\section{Levantamiento de información}
% =========================================================
\begin{frame}{Levantamiento de información}
  \begin{columns}[T]
    \begin{column}{0.32\textwidth}
      \begin{block}{1. Entrevista inicial}
        \small Plantilla de preguntas por servicio con el personal de la biblioteca.
      \end{block}
    \end{column}
    \begin{column}{0.32\textwidth}
      \begin{block}{2. Seguimiento}
        \small Preguntas de aclaración (16-10-2026): compras, carga de usuarios,
        multas, inventario y colecciones móviles.
      \end{block}
    \end{column}
    \begin{column}{0.32\textwidth}
      \begin{block}{3. Consolidación}
        \small Documento de levantamiento complementado, con correcciones y
        decisiones del equipo.
      \end{block}
    \end{column}
  \end{columns}

  \vspace{1em}
  \textbf{Resultado:} base común para construir los DFD y los BPMN, con
  registro explícito de supuestos y pendientes.
\end{frame}

% =========================================================
\section{Elementos identificados}
% =========================================================
\begin{frame}{Procesos identificados}
  \small
  \begin{tabularx}{\textwidth}{@{}l X@{}}
    \toprule
    \textbf{Proceso} & \textbf{Propósito} \\
    \midrule
    Acreditación          & Reportes estadísticos de uso para las carreras \\
    Catalogación          & Ingreso por donación o compra, catalogación y procesamiento técnico \\
    Circulación           & Préstamos, devoluciones, ingreso a sala, material digital y multas \\
    Servicios comunitarios& Actividades propias, de terceros y mandatos institucionales \\
    Inventario            & Verificar existencia física y depurar perdidos \\
    Descarte              & Retirar material según reglamento (meta: 1.000 copias/año) \\
    Colecciones móviles   & Préstamo anual de libros a departamentos \\
    Actualización usuarios & Carga masiva de usuarios y actualización de datos de usuarios \\
    \bottomrule
  \end{tabularx}
\end{frame}

\begin{frame}{Entidades externas}
  \begin{columns}[T]
    \begin{column}{0.5\textwidth}
      \begin{itemize}\small
        \item \textbf{Usuario} (estudiante, académico, funcionario)
        \item \textbf{Comunidad externa / Invitado}
        \item \textbf{Donante}
        \item \textbf{Jefe de carrera / Experto disciplinar}
        \item \textbf{Departamento académico}
      \end{itemize}
    \end{column}
    \begin{column}{0.5\textwidth}
      \begin{itemize}\small
        \item \textbf{Proveedor}
        \item \textbf{Institución receptora} (descarte)
        \item \textbf{Autoridades universitarias}
        \item \textbf{Contraloría}
        \item \textbf{DSI} (carga masiva de usuarios)
      \end{itemize}
    \end{column}
  \end{columns}
\end{frame}

\begin{frame}{Almacenes de datos}
  \small
  \begin{tabularx}{\textwidth}{@{}l l X@{}}
    \toprule
    \textbf{ID} & \textbf{Almacén} & \textbf{Contenido} \\
    \midrule
    A1 & Horizont                       & Catálogo, ítems, usuarios, préstamos, multas, estados \\
    A2 & Excel de registro de ingreso   & Registro manual de cada libro ingresado \\
    A3 & Drive de libros digitalizados  & Copias digitales (sin estadísticas de uso) \\
    A4 & Excel de inventario / perdidos & Resultado de cada inventario y perdidos acumulados \\
    A5 & Drive de actas de descarte     & Actas con los libros dados de baja \\
    A6 & Listado de colecciones móviles & Libros prestados a departamentos y su ubicación \\
    \bottomrule
  \end{tabularx}

  
\end{frame}

% =========================================================
\section{Diagramas de flujo de datos}
% =========================================================
\begin{frame}{DFD nivel 0 — Diagrama de contexto}
  \begin{columns}[c]
    \begin{column}{0.64\textwidth}
      \centering
      \figura{../Diagramas/DFD NIVEL 0.pdf}{8.5cm}{5.5cm}
    \end{column}
  \end{columns}
\end{frame}

\begin{frame}{DFD nivel 1}
  \begin{columns}[c]
    \begin{column}{0.64\textwidth}
      \centering
      \figura{../Diagramas/DFD NIVEL 1.pdf}{8.5cm}{5.5cm}
    \end{column}

  \end{columns}
\end{frame}

\begin{frame}{BPMN de Circulación}
  \centering
  \figura{../Diagramas/CIRCULACION.pdf}{13cm}{6.5cm}
\end{frame}


% =========================================================
\section{Hallazgos y próximos pasos}
% =========================================================
\begin{frame}{Hallazgos preliminares}
  \begin{itemize}
    \item \textbf{Reportes manuales:} acreditación y descarte dependen de filtrar
          Excel exportados de Horizont ``al ojo''.
    \item \textbf{Doble registro:} el Excel de ingreso y Horizont no están relacionados.
    \item \textbf{Sin trazabilidad digital:} el Drive no entrega estadísticas de uso.
    \item \textbf{Comunicación informal:} solicitudes y respuestas por correo, sin formato.
    \item \textbf{Ubicación física:} Horizont no la guarda; errores en estantería.
    \item \textbf{Ingreso a sala lento:} se registra como préstamo y devolución inmediata.
  \end{itemize}
\end{frame}

% =========================================================
\section{Propuesta de solución}
% =========================================================
\begin{frame}{Propuesta: NAS + BookOrbit}
  \begin{columns}[T]
    \begin{column}{0.52\textwidth}
      \begin{block}{Qué se propone}
        \small
        Reemplazar el almacén digital en Google Drive por una infraestructura local compuesta por un \textbf{NAS} y la plataforma \textbf{BookOrbit}.
      \end{block}
      \vspace{0.4em}
      \begin{itemize}\small
        \item Centraliza y respalda localmente la colección digital.
        \item Publica los documentos con acceso controlado mediante \textbf{SSO / OIDC} para cuentas institucionales.
        \item Vincula el catálogo \textbf{Horizont} con enlaces directos al lector digital.
      \end{itemize}
    \end{column}
    \begin{column}{0.44\textwidth}
      \begin{block}{Beneficios clave}
        \small
        \begin{itemize}
          \item Estadísticas de uso para acreditación.
          \item Mejor experiencia de lectura para estudiantes y docentes.
          \item Mayor control sobre seguridad, acceso y soberanía tecnológica.
        \end{itemize}
      \end{block}
    \end{column}
  \end{columns}
\end{frame}


% ---------------------------------------------------------

\end{document}
